\documentclass[../../../include/open-logic-section]{subfiles}

\begin{document}

\olfileid[tr]{sth}{card-arithmetic}{ch}

\olsection{Süreklilik Hipotezi}

Önceki sonuç, sonsuz kardinal sayıların $2$'nin kuvvetleriyle, yani
(\olref[opps]{lem:SizePowerset2Exp} gereği) kuvvet kümeleri almakla
``dolaysız biçimde davranacağı'' güvence altına alınsaydı kardinal üs
almanın oldukça \emph{kolay} olacağını (doğru biçimde) sezdirir. Fakat
sonsuz kardinal sayıların kuvvet kümeleriyle dolaysız biçimde
\emph{davrandığını} varsayamayız.

Bunu açıklamaya başlamak için kullanışlı bir gösterim sunuyoruz.

\begin{defn}
$\cardsucc{\cardfont{a}}$, $\cardfont{a}$'dan kesinlikle büyük olan en
küçük kardinal sayı olmak üzere iki sonsuz dizi tanımlarız:
\begin{align*}
	\aleph_{0} &\defis \omega &
	\beth_{0} &\defis \omega\\
	\aleph_{\alpha \ordplus 1} &\defis \cardsucc{(\aleph_{\alpha})} &
	\beth_{\alpha+1} &\defis \cardexpo{2}{\beth_{\alpha}}\\
	\aleph_{\alpha} &\defis \bigcup_{\beta< \alpha} \aleph_{\beta} &
	\beth_{\alpha} &\defis \bigcup_{\beta < \alpha}\beth_{\beta} &
	\text{$\alpha$ bir limit sıra sayısı olduğunda}.
	\end{align*}
\end{defn}

$\cardsucc{\cardfont{a}}$ tanımı yerindedir; çünkü
\olref[cardinals][classing]{lem:NoLargestCardinal}, her
$\cardfont{a}$ kardinal sayısı için $\cardfont{a}$'dan büyük bir kardinal
sayı bulunduğunu söyler ve Sonlu Ötesi Tümevarım, $\cardfont{a}$'dan
büyük \emph{en küçük} kardinal sayının bulunduğunu güvence altına alır.
Tanımın geri kalanı sonlu ötesi özyineleme ile sağlanır.

Cantor bu ``$\aleph$'' gösterimini ortaya koydu; bu, İbrani alfabesinin
ilk harfi ve İbranice ``sonsuz'' sözcüğünün ilk harfi olan \emph{aleph}'tir.
Peirce ``$\beth$'' gösterimini ortaya koydu; bu da İbrani alfabesinin
ikinci harfi olan \emph{beth}'tir.\footnote{Peirce bu gösterimi Aralık
1900'de Cantor'a yazdığı bir mektupta kullandı. Ne yazık ki aynı yerde
$\alpha\geq\omega$ için $\beth_\alpha$'nın var olmadığına ilişkin kötü
bir sav da ileri sürdü.} Bu gösterimler bize sonsuz kardinal sayılar verir.

\begin{prop}
Her $\alpha$ sıra sayısı için $\aleph_\alpha$ ve $\beth_\alpha$ kardinal
sayılardır.
\end{prop}

\begin{proof}
Her iki sonuç da basit bir sonlu ötesi tümevarımla çıkar.
$\aleph_0=\beth_0=\omega$, \olref[cardinals][classing]{omegaisacardinal}
gereği bir kardinal sayıdır. $\aleph_\alpha$ ve $\beth_\alpha$'nın
ikisini de kardinal sayı varsayarsak $\aleph_{\alpha+1}$ ve
$\beth_{\alpha+1}$ açıkça kardinal sayılar olarak tanımlanmıştır. Bir
kardinal sayılar kümesinin birleşimi de
\olref[cardinals][classing]{unioncardinalscardinal} gereği bir kardinal
sayıdır.
\end{proof}
\noindent
Üstelik her sonsuz kardinal sayı bir $\aleph$'tir:

\begin{prop}
$\cardfont{a}$ sonsuz bir kardinal sayıysa biricik bir $\gamma$ için
$\cardfont{a}=\aleph_\gamma$ olur.
\end{prop}

\begin{proof}
Kardinal sayılar üzerinde sonlu ötesi tümevarımla. Tümevarım için
$\cardfont{b}<\cardfont{a}$ ise
$\cardfont{b}=\aleph_{\gamma_\cardfont{b}}$ olduğunu varsayın.
Bir $\cardfont{b}$ için
$\cardfont{a}=\cardsucc{\cardfont{b}}$ ise
$\cardfont{a}=\cardsucc{(\aleph_{\gamma_\cardfont{b}})}=
\aleph_{\gamma_\cardfont{b}+1}$ olur. $\cardfont{a}$ hiçbir kardinal
sayının ardılı değilse, kardinal sayılar sıra sayıları olduğundan
$\cardfont{a}=\bigcup_{\cardfont{b}<\cardfont{a}}\cardfont{b}=
\bigcup_{\cardfont{b}<\cardfont{a}}{\aleph_{\gamma_\cardfont{b}}}$;
dolayısıyla
$\gamma=\bigcup_{\cardfont{b}<\cardfont{a}}\gamma_\cardfont{b}$ olmak
üzere $\cardfont{a}=\aleph_\gamma$ olur.
\end{proof}

Her sonsuz kardinal sayı bir $\aleph$ olduğuna göre şu soruyu sormak
doğaldır: her sonsuz kardinal sayı bir $\beth$ midir? Durum gerçekten
\emph{böyle} olsaydı sonsuz kardinal sayılar kuvvet kümesi alma işlemiyle
``dolaysız biçimde davranırdı''. Gerçekten de şunu elde ederdik:

\begin{defish}
\emph{Genelleştirilmiş Süreklilik Hipotezi} (GCH).
Her $\alpha$ için $\aleph_\alpha=\beth_\alpha$.
\end{defish}

Üstelik GCH geçerli olsaydı kardinal üs alma işleminde önemli
sadeleştirmeler yapabilirdik. Özellikle $\cardfont{b}<\cardfont{a}$
olduğunda $\cardexpo{\cardfont{a}}{\cardfont{b}}$ değerinin
$\cardfont{a}\leq\cardexpo{\cardfont{a}}{\cardfont{b}}\leq
\cardsucc{\cardfont{a}}$ arasında sıkıştığını gösterebilirdik. Daha sonra
iki olasılıktan hangisinin gerçekleştiğini (yani
$\cardfont{a}=\cardexpo{\cardfont{a}}{\cardfont{b}}$ mi, yoksa
$\cardexpo{\cardfont{a}}{\cardfont{b}}=\cardsucc{\cardfont{a}}$ mı
olduğunu) belirleyen kesin koşulları verebilirdik.\footnote{Koşulu
\emph{eşsonluluk} belirler.}

Fakat GCH bir \emph{hipotezdir}, \emph{teorem} değildir. Gerçekten
\citet{Godel1938}, $\ZFC$ tutarlıysa $\ZFC+\text{GCH}$'nin de tutarlı
olduğunu ispatladı. Daha sonra $\lnot$GCH'yi de $\ZFC$'ye ekleyebileceğimiz
ortaya çıktı. Gerçekten GCH'nin en basit ve önemsiz olmayan \emph{örneğini}
ele alın:

\begin{defish}
\emph{Süreklilik Hipotezi} (CH). $\aleph_1=\beth_1$.
\end{defish}

\citet{Cohen1963}, $\ZFC$ tutarlıysa $\ZFC+\lnot\text{CH}$'nin de
tutarlı olduğunu ispatladı. Dolayısıyla Süreklilik Hipotezi $\ZFC$'den
bağımsızdır.

``Süreklilik'', gerçek doğru $\Real$ için kullanılan başka bir ad olduğu
için bu hipoteze Süreklilik Hipotezi denir.
\olref[opps]{continuumis2aleph0}, $\card{\Real}=\beth_1$ olduğunu söyler.
Dolayısıyla Süreklilik Hipotezi, doğal sayıların kardinalitesi
$\aleph_0=\beth_0$ ile sürekliliğin kardinalitesi $\beth_1$ arasında
başka kardinal sayı bulunmadığını söyler.

(G)CH'nin $\ZFC$'den \emph{bağımsızlığı} karşısında bunların
\emph{doğruluğu} hakkında ne söylemeliyiz? Söylenecek \emph{çok} şey
vardır. Gerçekten küme teorisindeki verimli güncel çalışmaların çoğu bu
sorunları incelemeye yönelmiştir. Fakat iki kısa nokta kesinlikle
vurgulanmaya değerdir.

Birincisi: bu biçimsel bağımsızlık sonuçlarından GCH ya da CH'nin doğruluk
değeri bakımından \emph{belirsiz} olduğu \emph{doğrudan} çıkmaz. Belki de
yalnızca bize doğal görünen ve soruyu şu ya da bu yönde çözecek daha fazla
aksiyom eklememiz gerekir. Gödel'in kendisi doğru karşılığın bu olduğunu
öne sürmüştür.

İkincisi: CH'nin $\ZFC$'den bağımsız olması kesinlikle \emph{çarpıcıdır},
fakat (sözcüğün gerçek anlamıyla) \emph{inanılmaz} değildir. Mesele
yalnızca şudur: $\ZFC$'nin söylediği kadarıyla kardinal sayılardan onların
ardıllarına geçmek, yalnızca kuvvet kümeleri almaktan daha incelikli bir
araç gerektirebilir.
 %The operation of taking powersets moves us from one stage of the
 %hierarchy to its successor stage (from $V_{\alpha}$ to
 %$V_{\alpha+1}$).

Bu iki gözlemden sonra daha fazlasını öğrenmek istiyorsanız artık bu
tartışmada taraf olan çeşitli filozoflara ve matematikçilere yönelmeniz
gerekecek.\footnote{Başlangıç olarak
\citet[\S15.6]{Potter2004} okuyabilirsiniz.}

\end{document}
